% Complete root/B/D actual-file and exact primary-source audits PASS.
\begin{proposition}[Proved: a fixed canonical-defect annihilator and linear geometric genus]
\label{D:cover-uniform-canonical-defect}
Let $(R,\mathfrak m,k)$ be a complete equicharacteristic normal local
domain of dimension two, with algebraically closed coefficient field $k$.
There is a fixed $\mathfrak m$-primary ideal $J\subset R$ such that the
following holds. Let $A$ be a finite normal local domain over $R$, of
generic degree $n$, with $R\subset A$, and suppose that
\[
 \operatorname{Spec}A\setminus\{\mathfrak m_A\}
 \longrightarrow \operatorname{Spec}R\setminus\{\mathfrak m\}
\]
is etale. For a projective regular resolution
$\pi:Z\to\operatorname{Spec}A$ which is an isomorphism on the puncture,
use the natural trace inclusion $\pi_*\omega_Z\subset\omega_A$. Then
\[
 J\omega_A\subset\pi_*\omega_Z,
 \qquad JH^1(Z,\mathcal O_Z)=0.                         \tag{D.CD.1}
\]
If $Q=(x,y)\subset J$ is one fixed parameter ideal of $R$, then
\[
 p_g(A):=\dim_kH^1(Z,\mathcal O_Z)
 =\operatorname{length}_A(\omega_A/\pi_*\omega_Z)
 \le n\operatorname{length}_R(R/Q).                    \tag{D.CD.2}
\]
The constants and the annihilator are independent of the cover.
For a finite normal $R$-torsion-free algebra, the corresponding assertions
hold componentwise, and the sum of the geometric genera has the same bound
with $n$ the total generic degree. No separability restriction on $n$ is
imposed; etaleness on the puncture supplies generic separability.
\end{proposition}
\begin{proof}
\emph{A fixed different ideal.}
Choose generically separable Noether normalizations
$S_i=k[[x_i,y_i]]\subset R$, finite on $R$, such that the sum of their
Noether differents contains an $\mathfrak m$-primary ideal. Such a finite
choice exists: over the perfect field $k$, the Jacobian ideal is supported
on the isolated singular locus, is the sum of the Kahler differents of
Noether normalizations, and each Kahler different is contained in its
Noether different. A finite selection suffices by Noetherianity
\citep[Section 3, especially (3.1) and Examples 3.2--3.3,
pp.~6--7]{C_IyengarTakahashiStable}.
If $R$ is regular, its identity normalization can be used. Choose once an
$\mathfrak m$-primary ideal
\[
 J\subset\sum_i\mathfrak N(R/S_i).
\]
All rings $R$ and $A$ here are Cohen--Macaulay. In particular $R$ and $A$
are finite free over each $S_i$.

\emph{Trace duality on the punctured cover.}
Put $K=\operatorname{Frac}R$ and $L=\operatorname{Frac}A$.
Using the nondegenerate generic trace pairing, identify
\[
 \operatorname{Hom}_R(A,R)
 =\{c\in L:\operatorname{Tr}_{L/K}(cA)\subset R\}.
\]
This fractional $A$-module is exactly $A$. Indeed, at every height-one
prime of $R$ the finite algebra is etale and its trace dual is itself.
Every height-one prime of $A$ lies over such a prime of $R$, so an element
of the displayed trace dual belongs to every height-one localization of
$A$. Normality then puts it in $A$. Conversely the trace of an integral
element is integral over $R$ and lies in $K$, hence lies in $R$.
We have therefore proved, with its specified generic identification,
\[
 \operatorname{Hom}_R(A,R)=A.                          \tag{D.CD.3}
\]
Finite duality and transitivity identify
$\omega_A=\operatorname{Hom}_R(A,\omega_R)$.

\emph{A Noether different clears the canonical lattice.}
Fix $i$. Choose a generator of the dualizing module of the regular ring
$S_i$ and use finite duality to identify $\omega_R$ with
$\operatorname{Hom}_{S_i}(R,S_i)$ times that generator.
Write $\beta_i$ for the resulting trace section. It is nonzero generically.
The evaluation description of the Noether different is
\[
 \mathfrak N(R/S_i)=
 \operatorname{im}\bigl(\operatorname{Hom}_R(\omega_R,R)
       \longrightarrow R,\quad \varphi\longmapsto\varphi(\beta_i)\bigr),
                                                               \tag{D.CD.4}
\]
with the generator of $\omega_{S_i}$ suppressed in the notation.
This is Stacks Lemma 49.6.7, including its trace normalization
\citep[Tags 0BVT and 0DWM, Lemmas 49.6.7 and 49.15.1--2]
{C_StacksCanonicalDefect}.
The module $\omega_R$ is torsion-free of rank one. Consequently if
$z=\varphi(\beta_i)$, then $z\alpha=\varphi(\alpha)\beta_i$ for every
$\alpha\in\omega_R$, as one checks in its one-dimensional generic fibre.
It follows that
\[
 \mathfrak N(R/S_i)\omega_R\subset R\beta_i.
\]
Apply $\operatorname{Hom}_R(A,-)$ and use (D.CD.3). With the generic
trace identifications and trace transitivity, this gives
\[
 \mathfrak N(R/S_i)\omega_A\subset A\beta_i,            \tag{D.CD.5}
\]
where on the right $\beta_i$ now denotes the trace section for $A/S_i$.
The notation is compatible because
$\operatorname{Tr}_{L/\operatorname{Frac}S_i}
=\operatorname{Tr}_{K/\operatorname{Frac}S_i}\circ
 \operatorname{Tr}_{L/K}$.

\emph{Regularity of this section on the resolution.}
The composite $g_i:Z\to\operatorname{Spec}S_i$ is projective and
generically finite separable. It is an lci morphism of virtual relative
dimension zero. To see this without any finite-type-over-$k$ assumption,
embed $Z$ in a projective space over $S_i$; this is a regular immersion
because both the ambient scheme and $Z$ are regular. The fundamental
class of an lci morphism therefore gives
\[
 \mathcal O_Z\longrightarrow \omega_{Z/S_i}.
\]
After tensoring with the chosen generator of $\omega_{S_i}$, its image of
$1$ is a regular section of $\omega_Z$ whose generic value is precisely
$\beta_i$. Here the normalization is the canonical etale identification:
on the generic finite etale locus, relative duality identifies this image
with the trace functional. Equivalently, a local complete-intersection
presentation gives its Jacobian determinant times the Koszul dualizing
generator; on the etale locus that generator corresponds to the inverse
Jacobian under the trace pairing. These identifications are the ones in
Stacks Lemma 48.29.2 and the finite-flat comparison with duality
\citep[Tags 0E9Z and 0DWM; Lemmas 48.29.2 and 49.15.1--2]
{C_StacksCanonicalDefect}. Thus
$A\beta_i\subset\pi_*\omega_Z$. Summing (D.CD.5) proves the first
inclusion of (D.CD.1). In particular no use of ordinary absolute
differentials on a formal power-series field has been made.

\emph{Duality and length.}
The ring $A$ is complete excellent and has a dualizing complex.
Grauert--Riemenschneider vanishing for normal modifications of local
surfaces gives $R^1\pi_*\omega_Z=0$ in every characteristic
\citep[Tag 0AXD, Proposition 54.7.8, with the proof in Section 54.7]
{C_StacksCanonicalDefect}. We also have $\pi_*\mathcal O_Z=A$ and
$R^j\pi_*\mathcal O_Z=0$ for $j>1$. Put
$M=H^1(Z,\mathcal O_Z)$, a finite-length $A$-module. Duality applied to
$A\to R\pi_*\mathcal O_Z\to M[-1]$ gives
\[
 0\longrightarrow\pi_*\omega_Z\longrightarrow\omega_A
 \longrightarrow D_A(M)\longrightarrow0,               \tag{D.CD.6}
\]
where $D_A$ is Matlis duality. Thus $J$ kills $D_A(M)$ and hence $M$;
it also follows that the quotient in (D.CD.6) has length $p_g(A)$.
As an $R$-module, $\omega_A$ is maximal Cohen--Macaulay of rank $n$.
For any fixed parameter ideal $Q=(x,y)\subset J$,
\[
 \operatorname{length}_R(\omega_A/J\omega_A)
 \le\operatorname{length}_R(\omega_A/Q\omega_A)
 =n\operatorname{length}_R(R/Q).
\]
The last equality is the multiplicity formula for a maximal
Cohen--Macaulay module, or follows by taking ranks over $k[[x,y]]$,
over which both modules are finite free. Every residue field in the
connected finite cover is $k$, so these finite lengths are also the
stated $k$-dimensions. This proves (D.CD.2). A finite normal algebra over
a complete local ring is a finite product of normal local domains;
apply the argument to each dominating component and add their ranks.
\end{proof}

\begin{corollary}[Proved: cohomology of fixed pulled-back exceptional thickenings]
\label{D:cover-thickening-linear-cohomology}
In the algebraic or strict-henselian punctured-cover setting of
Proposition~\ref{D:punctured-cover-linear-resolution}, fix the positive
exceptional cycle $A_0=\sum_i a_iE_i$ on the base resolution $Y$.
For a degree-$n$ cover and the regular model $f:Z\to Y$ used there, put
$C=f^*A_0$. For each fixed integer $t\ge1$ there is a constant $c_t$,
depending only on the base germ, its fixed resolution, $A_0$, and $t$,
such that
\[
 h^0(\mathcal O_{tC})+h^1(\mathcal O_{tC})\le c_t n.
                                                               \tag{D.CD.7}
\]
The cohomology here is that of the full possibly nonreduced Cartier
cycle. The same assertion holds on the completed models.
\end{corollary}
\begin{proof}
Complete the base germ. Excellence makes the normalization and its
completion compatible and normal; finite covers split into complete
normal local factors with total generic degree $n$. Formal functions
and flat base change identify the finite-length $H^1$ of a resolution
with its completed version. Proposition~\ref{D:cover-uniform-canonical-defect}
therefore gives $h^1(Z,\mathcal O_Z)\le c_0 n$ with fixed $c_0$.
For every effective exceptional Cartier cycle $D$ the sequence
\[
 0\longrightarrow\mathcal O_Z(-D)\longrightarrow\mathcal O_Z
 \longrightarrow\mathcal O_D\longrightarrow0
\]
gives a surjection $H^1(Z,\mathcal O_Z)\to H^1(D,\mathcal O_D)$:
$H^2(Z,\mathcal O_Z(-D))=0$ because the proper map to the affine
local surface has fibres of dimension at most one. Hence
$h^1(\mathcal O_{tC})\le c_0n$.

The previously proved exceptional arithmetic calculation gives
\[
 C^2=nA_0^2,\qquad
 -n\sum_i a_i\le K_Z\cdot C\le nK_Y\cdot A_0
\]
after contracting the point-mapped rational $(-1)$-curves, as in
Propositions~\ref{D:exceptional-arithmetic-genus-bound} and
\ref{D:punctured-cover-linear-resolution}. The cycle adjunction formula
now gives
\[
 \chi(\mathcal O_{tC})
 =-\tfrac12\bigl(t^2C^2+tK_Z\cdot C\bigr)=O_t(n).
\]
Since $tC$ is proper of dimension one,
$h^0(\mathcal O_{tC})=\chi(\mathcal O_{tC})+h^1(\mathcal O_{tC})$.
This proves (D.CD.7). All assertions are insensitive to completing the
algebraic or strict-henselian germ, by the same proper base-change and
formal-functions argument.
\end{proof}

\begin{remark}[Scope of the remaining counting problem]
The first proposition applies to arbitrary complete equicharacteristic
normal local surface domains with algebraically closed coefficient field.
The second uses the previously specified algebraic or strict-henselian
resolution setup for its intersection estimates. These are uniform
annihilator and cohomology bounds. They do not provide a presentation of
cover-dependent module objects with $O(nr^2)$ parameters, and no hereditary
quadratic representation-count or finite-generation conclusion is inferred.
\end{remark}
